% Content for the organizers' Overleaf template (acmart, sigconf). Paste each block over the
% gray \guide{...} text of the matching place in main.tex; leave the black text as it is.
% Upload report/figures/pipeline.pdf and report/figures/examples.pdf to the project.
% Author block for the template:
%   \author{Islam Kossymbayev}
%   \affiliation{\institution{Nazarbayev University}\city{Astana}\country{Kazakhstan}}
%   \email{islam.kossymbayev@gmail.com}
% Check the city before you compile.

%%%%%%%%%%%%%%%%%%%% TITLE (replace \guide{Your Title Here}) %%%%%%%%%%%%%%%%%%%%
% Disk-Scale ConvNeXt and Crop-Level ResNet U-Nets with Exact Candidate Selection

%%%%%%%%%%%%%%%%%%%% ABSTRACT (after the black sentence) %%%%%%%%%%%%%%%%%%%%
We predict a filament probability map with two groups of U-Nets, twelve ResNet34 models that
see 512-pixel crops and three ConvNeXt-Tiny models that see the whole solar disk at 1536 pixels,
and average their flip-augmented outputs. Instances are not cut at one threshold. We cut the map
at six thresholds, which yields a forest of nested candidate masks, score each candidate with a
LightGBM regressor that predicts its matched IoU, and choose the best set of disjoint candidates
exactly by dynamic programming. The score is exact because disjoint candidates make Panoptic
Quality additive over candidates. On 707 training images, scored out of fold with date-grouped
folds, this reaches PQ 0.443, against 0.379 for a single ResNet34 with the plain hysteresis rule.
The public leaderboard score is 0.38.

%%%%%%%%%%%%%%%%%%%% 1 INTRODUCTION (replace the gray sentence) %%%%%%%%%%%%%%%%%%%%
Our entry has three parts: a probability map from an ensemble of U-Nets
(Section~\ref{sec:method}), a candidate forest with exact selection that targets the metric
directly, and a validation protocol that keeps every number out of fold
(Section~\ref{sec:eval}). We only used the competition training annotations. Test images were
never labelled and never used for training, and no external filament annotations were used.
The whole-disk ConvNeXt model follows ideas of a public entry
(github.com/pragyaangaur/NSO-Solar-Filament-Challenge); the code was written from scratch.
The encoders are ImageNet-pretrained timm and torchvision weights.

%%%%%%%%%%%%%%%%%%%% 2 METHODOLOGY (replace the gray list) %%%%%%%%%%%%%%%%%%%%
\label{sec:method}
Figure~\ref{fig:pipeline} shows the pipeline.

\begin{figure}[t]
\centering
\includegraphics[width=\linewidth]{pipeline.pdf}
\caption{End-to-end pipeline for one test image.}
\label{fig:pipeline}
\end{figure}

\paragraph{Data and splits.}
The 707 training images (1154 annotator views, 8199 filaments) are split into three folds
grouped by calendar date, so same-day images never cross folds. Every development number below
is out of fold: an image is scored only by models trained on the other two folds. Test maps use
all models. A training sample uses one random annotator view of its image.

\paragraph{Crop-level ResNet34 U-Net.}
Input channels are intensity and a local contrast response (Gaussian background, $\sigma=20$,
minus the image, scaled to $[-1,1]$). Training uses 512-pixel crops centred near annotated
filaments, flips, BCE with positive weight 10 plus soft Dice, AdamW, and 30 epochs. Each fold has
four members (two seeds, two view-sampling rules), twelve models in total. Inference uses 512-pixel
tiles with Hann blending and four flip augmentations.

\paragraph{Whole-disk ConvNeXt U-Net.}
The full disk is resized to 1024 pixels for 30 epochs and then to 1536 pixels for 24 epochs
(ConvNeXt-Tiny encoder, five-block decoder). Besides the filament output, a second output
predicts instance boundaries (weight 0.2). We use flips and rotations by 90$^\circ$, photometric
and noise augmentation, AdamW with warmup and cosine decay, and an EMA of the weights. One model per
fold gives three models. Flip-averaged predictions are resized back to 2048 pixels.

\paragraph{Candidate forest and exact selection.}
Let $p$ be the average of the ResNet group and the ConvNeXt group. For each grow threshold in
$\{0.3,\dots,0.8\}$ we close the mask with a $3\times3$ ellipse, keep 8-connected components of
at least 30 pixels inside the disk minus a 65-pixel limb band, and erode each by one pixel.
Thresholding and closing are increasing, so each component lies inside exactly one component of
the next lower threshold, and the candidates of an image form a forest. Two candidates that are
not nested stay disjoint, hence any set without nested pairs (an antichain) is a valid
non-overlapping prediction, and a ground-truth filament exceeds IoU 0.5 with at most one of its
members.

Then PQ is additive. Keeping a candidate seen by $n$ annotator views adds its matched IoU sum
to the numerator and $n/2$ to the denominator. With target PQ $=2t$, the best antichain maximises
$\sum (\hat s - t)$ over kept candidates, where $\hat s$ is the candidate's mean matched IoU per
view. A LightGBM regressor predicts $\hat s$ from 24 features (probability statistics, seed
sizes, shape, contrast against a surrounding ring, radial position, observatory, and four tree
features: level, area over parent, number of children, largest child fraction). A dynamic program
over each forest returns the exact maximiser. The threshold $t$ is chosen by date-grouped inner
cross-validation on the training folds only. The kept masks are written as COCO RLE.

\begin{figure*}[t]
\centering
\includegraphics[width=0.78\textwidth]{examples.pdf}
\caption{Three test images (no labels exist for them): blended probability (top) and the
selected instances of the submitted file (bottom).}
\label{fig:examples}
\end{figure*}

%%%%%%%%%%%%%%%%%%%% 3 EVALUATION (replace the gray list) %%%%%%%%%%%%%%%%%%%%
\label{sec:eval}
\paragraph{Strategy.}
PQ follows the organizers' notebook (match at IoU $>0.5$, every annotator view scored against
the same prediction, counts pooled). Intervals are paired 95\% bootstrap intervals over
calendar dates (10{,}000 replicates). All choices were fixed on the development folds, and the
threshold is learned inside the training folds.

\begin{table}[t]
\caption{Out-of-fold PQ on 707 training images. $\Delta$ is the gain over the row above, with a
paired date-bootstrap 95\% interval. Public PQ is given where a file was submitted.}
\label{tab:results}
\small
\begin{tabular}{@{}p{3.6cm}rrr@{}}
\toprule
Method & PQ & $\Delta$ [95\% CI] & Public \\
\midrule
ResNet34, single, rule & 0.379 & & 0.32 \\
12 ResNet34, flip, rule & 0.396 & & \\
+ LightGBM scorer & 0.422 & +0.027 [+0.022, +0.031] & 0.36$^{\ast}$ \\
ConvNeXt 1536 + scorer & 0.425 & & \\
ResNet + ConvNeXt, scorer & 0.440 & +0.017 [+0.013, +0.022] & \\
+ multi-threshold forest & 0.443 & +0.005 [+0.002, +0.008] & 0.38 \\
\midrule
+ second ConvNeXt (negative) & 0.441 & $-$0.003 & \\
\bottomrule
\end{tabular}

\smallskip
{\footnotesize $^{\ast}$ six-model ensemble. The first row is the mean of members, no augmentation.
The second row is for four-member fold ensembles. The forest result varied between 0.4434 and
0.4443 over two runs of the same code (LightGBM nondeterminism); the gain over the single
threshold is positive in all three folds.}
\end{table}

\paragraph{Results.}
Table~\ref{tab:results} lists the steps. The largest gain comes from replacing the hand rule with
the scorer, which recovers filaments whose seed falls below the old 0.98 requirement. The disk-scale
ConvNeXt is as good alone as the whole ResNet ensemble with the scorer, and the two are
complementary. The forest raises the ceiling of selection (true-IoU oracle) from about 0.50 to 0.55,
but the scorer captures only a small part of that, so the net gain is small. A second ConvNeXt
(different pretraining, 1024 px) did not help in any selection setting.

\paragraph{Strengths and weaknesses.}
The selection is exact for the metric, adds no hand-tuned size or seed rule, and its threshold is
learned without touching held-out folds. Public PQ has stayed near 0.85 times development PQ so far (0.32 against 0.379, 0.36 against 0.419, 0.38 against 0.443);
we have not isolated the cause. Weaknesses: the scorer, not the segmentation
network, is now the bottleneck (oracle 0.55 against 0.443 achieved). Thin, faint filaments lack a
reliable candidate, and merges of touching filaments are only fixed when a higher threshold splits
them. Further ideas: instance-aware losses that train on candidate quality directly, a scorer with
image-level context, and more diverse encoders.
